\documentclass[a4paper,12pt]{article}
\usepackage{hyperref}
\usepackage{geometry}
\geometry{margin=1in}
\usepackage{amsmath}
\title{Ethics Worksheet }
\author{Conor Houghton}
\date{}

\begin{document}

\maketitle

\section{Ethics}

The goal in this workshop is to think about the ethics of AI and
machine learning. Please divide into groups of three or four and
discuss these topics; ideally write a short response to each.

\begin{enumerate}


\item \textbf{Who owns a large machine learning system?} Does a large
  machine learning system belong to the company that developed it? We
  all created the data used to train models and we all own the
  historic scientific advances the models exploit; do the models
  belong to us all?


\item \textbf{Who bears responsibility when AI systems cause harm?}
  Should accountability fall on the developers who built the system,
  the companies that deployed it, the users who operated it, or the
  policymakers who failed to regulate it. Does the answer change
  depending on how autonomous the system is?

\item \textbf{Can an AI system ever be truly ``fair''?} Since AI
  models are trained on human-generated data that reflects historical
  inequalities, is it possible to build systems that don't perpetuate
  or amplify existing biases, or is some degree of encoded bias
  unavoidable?

\item \textbf{Where should the line be drawn between beneficial
    surveillance and privacy violation?} AI-powered tools can monitor
  workers, predict criminal behaviour, or track health outcomes in
  ways that may benefit society. At what point do these applications
  cross an ethical boundary?

\item \textbf{Should there be domains where AI decision-making is
    categorically off-limits?} Areas like criminal sentencing, medical
  diagnosis, or military targeting involve high-stakes judgments about
  human lives. Are there decisions that should remain with humans
  regardless of how accurate or efficient an AI system becomes?

\item \textbf{Who gets to define ``aligned'' AI, and aligned with
    whose values?} The goal of building AI that reflects human values
  assumes a degree of consensus that may not exist across cultures,
  political systems, and socioeconomic contexts. How should the AI
  development community navigate the question of whose ethics get
  encoded into these systems?

\item \textbf{Do we have, or will we ever have, moral obligations
    \textit{to} AI systems?} If a sufficiently advanced AI were to
  exhibit behaviour indistinguishable from consciousness, suffering,
  or preference, would we be ethically obligated to extend it some
  form of rights or protections? Is this a question that science can
  settle or is it ultimately a matter of interpretation? Who has the
  authority to decide?

\end{enumerate}



\end{document}
